
The simulation of Section~\ref{sec:sim} reported, in an earlier version of
this round, that every interval construction fails on the noisy world:
coverage \coverageNoisyBefore\ against a nominal $95\%$, with the estimator
sitting \biasNoisyBefore\ below $V_{\mathrm{limit}}$ and a mean interval width
of \widthNoisyBefore. A gap the size of the whole interval is not something a
resampling scheme repairs, and the natural write-up was that the estimator is
biased in that world and the interval inherits the bias.

That write-up would have been wrong, and the two experiments below are what
stopped it. Both were designed so that either could refute the bias reading,
and both did.

\paragraph{The gap against sample size}
A finite-sample bias closes as $n$ grows. The increment was estimated at
\nSizeGrid\ sample sizes from \sizeLo\ to \sizeHi\ in every world, at
\nScaleReps\ replicates each, with the penalty held fixed. The sparse world
behaves as a finite-sample gap must: \sparseGapSmall\ at $n = 4{,}000$ falling
to \sparseGapLarge\ at \sizeHi. The noisy world does not: \noisyGapSmall\ at
$n = 4{,}000$ and still \noisyGapLarge\ at \sizeHi. The sparse world is the
control that makes the noisy result mean something --- without it, ``the gap
did not close'' could be a statement about the experiment rather than about
the world.

\paragraph{The gap against the penalty}
A fixed $L_2$ penalty is not scale-free, so it can masquerade as bias: it
shrinks coefficients materially at $n = 4{,}000$ and barely at $n = 120{,}000$,
which would make a large-sample fit a different specification rather than the
same one at its limit. Holding $n$ at $4{,}000$ and weakening the penalty by
\nPenaltyDecades\ orders of magnitude moves the noisy world's gap from
\noisyGapTight\ to \noisyGapLoose\ --- that is, not at all.

\paragraph{What was actually wrong}
Neither explanation surviving is what sent us to examine the target rather
than the estimator, and the target was wrong. The simulation enumerated the
generator's cells, carrying the \emph{true} register value, while in the noisy
world the estimator only ever sees a value that is replaced by a uniform draw
with probability \simNoiseRate. Both estimands were therefore computed on a
population the estimator does not sample from. Marginalising the noise exactly
(Section~\ref{sec:sim}) moves $V_{\mathrm{limit}}$ from \limitNoisyBefore\ to
\limitNoisyAfter. The estimator's mean was \meanNoisyEstimate\ throughout.
There was no bias; there was a wrong estimand in our own simulation, and it
took two arms of evidence against everything else to find it.

\paragraph{The same two experiments, after the correction}
Both arms were re-run against the corrected target, and there is nothing left
for them to explain: the noisy world's gap is \noisyGapFixed\ at $n = 4{,}000$
and \noisyGapFixedLarge\ at \sizeHi. The sparse world's gap is unchanged in
both runs, because the correction touches only the noisy world, which is what
makes it a control rather than a coincidence. Coverage on the noisy world is
\coverageNoisyAfter. \texttt{scripts/s18\_bias\_scaling.py -{}-legacy-target}
regenerates the pre-correction figures above, so the evidence for this
appendix is reproducible rather than quoted from a log we happen to have kept.

The general lesson is the one this paper makes elsewhere about the design
space: a quantity is defined relative to a declared population, and a
simulation's population is a declaration like any other. A quality mechanism
that acts on the data --- which is exactly what the register-quality axis of
Section~\ref{sec:surface} does on real logs --- changes what the estimand is,
not merely how well it is estimated. The register-quality results of
Section~\ref{sec:quality} are reported against the degraded register for that
reason, and the correction register (Appendix~\ref{app:corrections}) records
this as the round's largest error of its class.
